\scanpage{46}

% B1 | text | cancelled opening, with addition above it
\deleted{sont \uline{\unsure{canoniquement}} équivalentes, \inserted{l'équivalence devenant \unsure{à|ét.} isom près} Si on \unsure{veut} les deux
essentiellement identifier entre eux les différents $(C,\alpha)$,
on trouve ainsi une catégorie de \uline{Deuring}
\uline{canonique}}

% B2 | text | upper paragraph, continuation after cancellation
sont équivalents : la catégorie des équivalences
de $(C,\alpha)$ avec $(C',\alpha')$ est un groupoïde connexe de
$\pi_1$ $\mathbb Z/2\mathbb Z$ : i.e. \deleted{\unsure{de}} \unsure{toutes|deux} les équivalences sont
isom, et si on a deux équiv. $\varphi,\varphi'$, il y a \uline{à un
signe près} un seul isom. entre $\varphi$ et $\varphi'$.

% B3 | text | main middle paragraph
Si on se donne maintenant un couple $(C,\alpha)$,
il lui est associé une \deleted{torseur sous} gerbe sur $\mathbb A=\prod_{\substack{\ell\ (\text{y compris}\\\ell=p\text{ et }\ell=\infty)}}\mathbb Z_\ell$
\inserted{de lien $\mu_2$,} savoir celle des équivalences \inserted{de $(C,T)$} avec le \deleted{\unsure{gerbe}} couple
type de $(C(\mathbb A),T)$, qui en la place $p$ soit
compatible avec la donnée de $\alpha$. Cette gerbe
$\unsure{\mathscr G|\mathfrak G|\mathcal G}_{(C,\alpha)}$ est déterminée, ``à un signe près'' \deleted{\unsure{ou}} équivalence
unique près, indépendamment du choix de $(C,\alpha)$
(puisqu'il en est ainsi de $(C,\alpha)$ lui-$=$\ldots), \deleted{et}
$\deleted{\hat{\mathbb Z}=\mathrm{Gal}(\bar{\mathbb F}_p/\mathbb F_p)}$ \deleted{y opère trivialmt.} mais elle dépend
a priori, \inserted{$=$ pour $(C,\alpha)$ fixé,} du choix d'un $\bar{\mathbb F}_p$ \deleted{(ou du moins
de $\mathbb F_{p^2}$)} -- tout comme la structure $\unsure{-|\alpha}$ elle-$=$ !

% B4 | text | lower left, struck marginal note
% slanted
\margin{\deleted{car le $F$-cristal de rang $2$, pente $\frac12$ sur $\bar{\mathbb F}_p$ se \unsure{descendait|descendant} \uline{canoniquement} à $\mathbb F_{p^2}$}}

% B5 | text | last paragraph
On peut s'en tirer en définissant \uline{canoniqt} un
$\mathbb F_{p^2}$ (ou $=$ un $\bar{\mathbb F}_p$, si on y tenait\ldots) par les \unsure{polynomes|polynômes}
cyclotomiques $\Phi_{p^n-1}$,

% B6 | display | foot
\[
\bar{\mathbb F}_p=\varinjlim_n\left(\mathbb F_{p^n}=\mathbb F_p[T]/\Phi_{p^n-1}(T)\right)
\qquad \ldots\ldots
\]
\unsure{;}
% archive sheet number: 45
